%!TeX root = main.tex
%!encoding = UTF-8 Unicode
\section{Dynamic Measurements for Power Electronics}

In the field of power electronics, dynamic measurements are crucial for understanding the behavior of power electronic devices and systems under various operating conditions. 
These measurements help in characterizing the performance, stability, and efficiency of power converters, inverters, and other related components.

Table \ref{tab:dynamic_measurements} summarizes the key dynamic measurements commonly performed in power electronics, which are explained in detail in the following subsections.

\begin{table}[htbp]
\caption{Dynamic measurements and verification criteria}
\label{tab:dynamic_measurements}
{%
\renewcommand{\tabularxcolumn}[1]{m{#1}}
    \begin{tabularx}{\textwidth}{>{\raggedright\arraybackslash}X>{\raggedright\arraybackslash}X>{\raggedright\arraybackslash}X>{\raggedright\arraybackslash}X}
        \toprule
        \textbf{Test} & \textbf{Verification test} & \textbf{Problems} & \textbf{Verification} \\
        \midrule
        Loading & Impedance measurement & Influencing DUT operation & 1 decade higher/lower \\
        \midrule
        Small signal response & VNA measurements & - Bandwidth \newline - gain flatness & $f_{\text{3dB}} > f_{\text{interest,max}}$ \\
        \midrule
        Large signal response & Pulse generator & Nonlinearities: \newline - Saturation \newline - Hysteresis \newline - ... & $t_r > t_{\text{r,meas,min}}$ \\
        \midrule
        Parasitic coupling & Measurement in parallel to verified sensor in setup & - Capacitive coupling of voltage transient \newline - Inductive coupling of current transient & Matching with verifying sensor \\
        \midrule
        Transmission line & Measurement in parallel to verified sensor in setup & - Wave reflection \newline - Signal distortion & Matching with verifying sensor \\
        \bottomrule
    \end{tabularx}
}
\end{table}



\subsection{Probe Loading and Circuit Interaction}

\subsection{Small-Signal Response and Bandwidth Criteria}

\subsection{Large-Signal Response and Nonlinearities}

\subsection{Parasitic Coupling Paths}

\subsection{Transmission-Line Effects}

